\documentclass[12pt]{article}

% --- Packages ---
\usepackage[utf8]{inputenc}
\usepackage[T1]{fontenc}
\usepackage{geometry}
\geometry{margin=2cm}
\usepackage{hyperref}
\usepackage{enumitem}
\usepackage{booktabs}
\usepackage{titlesec}

% --- Custom Formatting ---
\titleformat{\section}{\large\bfseries}{\thesection}{1em}{}
\hypersetup{colorlinks=true, linkcolor=blue, urlcolor=cyan}

% --- Document Info ---
\title{Reading Report: Emergent Digital Technologies}
\author{Your Name}
\date{\today}

\begin{document}

	\maketitle

	\section{Article Information}
	\begin{description}[style=multiline, leftmargin=3cm, font=\bfseries]
		\item[Title:] the paper's title
		\item[Authors:] the paper's authors
		\item[Journal:] the paper's journal
		\item[Year:] the paper's year
		\item[DOI:] \href{https://doi.org/}{the paper's DOI}
	\end{description}

	\section{Summary of the Technology}
		Write a brief overview of the paper here...

	\section{Core Innovations}
		\begin{itemize}
			\item \textbf{Innovation 1:} ...
			\item \textbf{Innovation 2:} ...
			\item \textbf{Innovation 3:} ...
		\end{itemize}

	\section{Key Results and Performance}
		Summarize the findings.  e.g., In the 90\% accuracy in the New England test site the authors tested the algorithm in WRS Path 12, Row 31, achieving:

		\begin{itemize}
			\item \textbf{Something:}
			\item \textbf{Something else:}
		\end{itemize}

	\section{Critique and Emergent Potential}
		Why is this important?  i.e. This technology is foundational for modern cloud-based geospatial applications like ...

	\section{Discussion Questions}
		Formulate and discusse broad questions such as:
		\begin{enumerate}
			\item Could this algorithm be adapted for higher-resolution data like Sentinel-2?
		\end{enumerate}

\end{document}
